\uuid{F91t}
\titre{Absence de mémoire de la loi exponentielle}
\auteur{Maxime Nguyen}
\organisation{AMSCC}
\datecreate{2026-11-24}
\niveau{L2}
\module{Probabilité}
\chapitre{Probabilité continue}
\sousChapitre{Loi exponentielle}
\difficulte{2}
\theme{loi exponentielle, absence de mémoire, probabilité conditionnelle}

\contenu{
\begin{question}
Soit $\lambda > 0$ et $X$ une variable de loi exponentielle de paramètre $\lambda$, notée $X \sim \mathcal{E}(\lambda)$ : pour tout $t \geq 0$, $\mathbb{P}(X > t) = e^{-\lambda t}$. La probabilité conditionnelle $\mathbb{P}(X > 5 \mid X > 3)$ est égale à :
\begin{answers}
\good{$\mathbb{P}(X > 2)$}
\bad{$\mathbb{P}(X > 5)$}
\good{$e^{-2\lambda}$}
\bad{$1 - e^{-2\lambda}$}
\good{$\dfrac{\mathbb{P}(X > 5)}{\mathbb{P}(X > 3)}$}
\bad{$\mathbb{P}(X > 8)$}
\end{answers}
\begin{explanations}
E est la définition de la probabilité
conditionnelle (car $\{X > 5\} \subset \{X > 3\}$) ; elle vaut
$\frac{e^{-5\lambda}}{e^{-3\lambda}} = e^{-2\lambda} = \mathbb{P}(X > 2)$.

Contre-sens visés : croire que l'information « $X > 3$ » ne change rien (B) ;
confondre $\mathbb{P}(X > t)$ et $\mathbb{P}(X \leq t)$ (D) ; appliquer
l'absence de mémoire en ajoutant les durées au lieu de les retrancher (F).
\end{explanations}
\end{question}
}
